\section{De la guía con CLIP a la guía con modelos de difusión}

\subsection{Generación 3D guiada por CLIP}

Antes del surgimiento de Score Distillation, los investigadores exploraron primero el uso de conocimientos previos del modelo CLIP (Contrastive Language-Image Pre-training) para asistir en la reconstrucción 3D.

\begin{knowledgebox}{Introducción al modelo CLIP}
CLIP es un modelo de alineación visual-lingüística propuesto por OpenAI, que mapea imágenes y textos al mismo espacio de embeddings. Al calcular la distancia entre los vectores de embeddings, se puede medir el grado de alineación semántica entre una imagen y una descripción textual. CLIP se entrena en grandes conjuntos de datos de pares imagen-texto y posee una potente capacidad de transferencia sin ejemplos (zero-shot).
\end{knowledgebox}

\subsubsection{Reconstrucción 3D con pocos ejemplos}

La estrategia que utiliza la idea de CLIP es la siguiente:
\begin{itemize}
    \item Para las vistas que tienen imágenes de referencia: se utiliza la pérdida de reconstrucción L2 tradicional.
    \item Para las vistas que no tienen imágenes de referencia: se utiliza la distancia CLIP como función de pérdida.
\end{itemize}

Por ejemplo, al reconstruir un modelo de bulldozer (推土机), sabemos que, sin importar desde qué ángulo se observe, debería ser siempre ``bulldozer''. Por lo tanto, para las vistas que carecen de imágenes de referencia, podemos exigir que la imagen renderizada esté tan cerca como sea posible del texto ``bulldozer'' en el espacio de embeddings de CLIP.

\subsubsection{Generación 3D sin ejemplos}

El caso límite de la generación 3D guiada por CLIP es: \textbf{completamente sin ninguna imagen de referencia}, generando un modelo 3D basándose únicamente en una descripción textual. En este escenario, todas las vistas utilizan la pérdida CLIP:

$$
\theta^* = \arg\min_{\theta} \sum_{\pi \sim \mathcal{P}} d_{\text{CLIP}}(g(\theta, \pi),\ \text{``a bulldozer''})
$$

\begin{importantbox}{Cambio de paradigma: de la reconstrucción a la generación}
Cuando el número de imágenes de entrada disminuye hasta cero, la naturaleza de la tarea cambia fundamentalmente: deja de ser ``reconstrucción 3D'' para convertirse en ``generación 3D''. Esto representa el primer intento importante de crear contenido 3D utilizando conocimientos previos de 2D.

Trabajo representativo: DreamFields (2021) — logró la generación texto a 3D utilizando únicamente la pérdida CLIP y NeRF.
\end{importantbox}

\subsection{Limitaciones de la guía con CLIP}

Aunque la generación 3D guiada por CLIP abrió un nuevo camino, la calidad de salida es limitada; los modelos 3D generados suelen ser bastante toscos y carecen de detalles. Esto se debe a que CLIP es esencialmente un modelo discriminador (discriminator), que solo puede determinar si una imagen coincide con un texto, pero no puede ofrecer una guía detallada sobre ``cómo debería verse la imagen''.

\begin{warningbox}{Limitaciones de CLIP como función de pérdida}
La distancia en el espacio de embeddings de CLIP solo puede medir la alineación a nivel semántico y no puede proporcionar una guía de generación a nivel de píxel. Al optimizar con la pérdida CLIP, el modelo podría producir resultados ``semánticamente correctos pero visualmente poco naturales'' (problema de adversarial examples).
\end{warningbox}

\subsection{Pregunta central: ¿Podemos reemplazar CLIP con modelos de difusión?}

CLIP es un discriminador natural —ingresa una imagen y un texto, y devuelve una puntuación de coincidencia. Sin embargo, los modelos de difusión son generadores —ingresan un texto y producen una imagen.

\textbf{Pregunta}: ¿Cómo se puede utilizar un modelo generador como discriminador?

Si se utilizara GAN, esto sería sencillo —el propio GAN incluye un discriminador—. Pero los modelos de difusión no tienen un discriminador explícito. Score Distillation Sampling (SDS) fue propuesto precisamente para resolver este problema.

\subsection{Resumen de la sección}

La generación 3D guiada por CLIP (como DreamFields) demostró la viabilidad de generar contenido 3D utilizando conocimientos previos de 2D, pero está limitada porque CLIP solo puede ofrecer una guía a nivel semántico. La idea siguiente más natural es: reemplazar CLIP con modelos de difusión que posean un conocimiento más rico a nivel de píxel, logrando así resultados de generación de mayor calidad.

\newpage
